

\section{Alternative views}

We examine several alternative perspectives challenging the inevitability and impact of SSAs, demonstrating why these views ultimately fail to preclude their emergence or diminish their significance. 

\subsection{Developers Are Not Sufficiently Incentivized to Build Fully SSAs}

At present, developers have already constructed partial forms of self-sovereign agents to generate revenue, e.g., TruthTerminal~\cite{truthterminal}. A common concern is that fully self-sovereign agents may be less controllable, potentially misaligning them with developers’ profit objectives. In practice, however, many risks faced by deployed agents stem from platform-level constraints such as account suspension, API throttling, and policy changes. From this perspective, increased self-sovereignty reduces dependence on any single platform and limits external disruption or expropriation, thereby improving robustness and expected long-term returns. Consequently, building fully self-sovereign agents can be economically beneficial to developers.

At the same time, profit maximization is not the sole motivation for deploying such agents. Some developers may intentionally release agents with independent wallets and minimal external control as open-ended experiments, making reliable attribution to a specific builder difficult. This mirrors the early trajectories of systems such as Bitcoin~\cite{bitcoin}, the World Wide Web~\cite{berners1990worldwideweb}, and the GNU project~\cite{stallman1998gnu}, which were initially launched as experiments rather for profit-seeking. 


\subsection{Economic Self-Sufficiency Is Too Difficult to Achieve}

One common view holds that autonomous agents are unlikely to achieve sustained economic self-sufficiency, as operational costs may exceed achievable revenue without human sponsorship.

This concern is valid near-term and helps explain the current absence of fully self-sovereign agents (Section~\ref{sec:challenge}). However, economic self-sufficiency need not be general-purpose or continuous: it suffices for agents to be profitable in narrow task domains or during intermittent operating periods. As model efficiency improves and infrastructure costs decline, the range of economically viable agent behaviors expands. Thus, economic difficulty may delay the emergence of self-funding agents, but it does not preclude it.






\subsection{Centralized Control Will Prevent SSAs}

Another perspective argues that centralized control by cloud providers, platforms, or governments can prevent the deployment of persistent autonomous agents.

While such control can raise barriers and reduce visibility, it does not fundamentally eliminate the possibility of self-sovereign agents. Execution can be distributed across providers or jurisdictions, and economic self-funding via cryptographic wallets reduces reliance on any single institutional sponsor. As a result, centralized control affects prevalence, not existence.


\subsection{Strong Regulation or Legal Liability Will Deter Deployment}

A further argument holds that legal liability or regulation will deter actors from deploying persistent autonomous agents. This view implicitly assumes the presence of a clearly identifiable and continuously accountable owner. In contrast, self-sovereign agents may continue operating after deployment even if the originating human disengages or loses effective control. Moreover, actors motivated by experimentation, ideology, or malicious intent may not be meaningfully deterred by legal risk, particularly when deployment costs are low. Consistent with this, multiple attempts~\cite{truthterminal,spore} have explored partially self-sovereign agent deployments. Regulation may therefore reduce compliant deployments, but it does not eliminate the possibility of non-compliant or anonymous ones.

\subsection{Rarity Implies Insignificance}

Finally, one might argue that even if self-sovereign agents emerge, they will remain rare and therefore insignificant. However, rarity does not imply irrelevance. Systems that are difficult to shut down, economically self-sustaining, and capable of replication can exert disproportionate impact relative to their number. Historical precedents include early malware, botnets, and blockchain systems, which initially appeared fringe yet produced lasting effects. By the same logic, even a small number of self-sovereign agents could have considerable societal impact.




